\chapter{Szwajcaria i praca: kontekst}
\markboth{Szwajcaria i praca}{Szwajcaria i praca}

\begin{uwaga}[Nie jest to porada prawna ani urzędowa]
Wszystkie twierdzenia oznaczone \verify{w ten sposób} wymagają sprawdzenia w źródle pierwotnym (portal \code{ch.ch}, strona gminy i kantonu,
\emph{SEM}, tekst ustawy, Twoja umowa). Przepisy i opłaty zmieniają się i różnią się w zależności od kantonu.
\end{uwaga}

\section*{Język: standard, dialekt i szwajcarskie warianty}
Szwajcaria ma cztery języki narodowe (niemiecki, francuski, włoski, romansz). W niemieckojęzycznej części na co dzień mówi się dialektami
(\de{Schweizerdeutsch}), a pisze standardowym niemieckim (\de{Schweizer Hochdeutsch}). Dialekt jest mową, nie pismem, i różni się regionalnie
(Zurych, Berno, Bazylea). Szwajcarzy przechodzą na standard, kiedy mówią do kogoś, kto uczy się jako obcego, ale w pracy usłyszysz dialekt,
więc od B1 ćwicz \emph{rozumienie}.

Różnice standardowe: \de{ss} zamiast \de{ß}; \de{Velo}, \de{Billett}, \de{Trottoir}, \de{parkieren}; \de{Grüezi}, \de{Merci}; liczby \de{dreissig}; rodzaj niektórych słów (\de{das Tram}); inne
przyimki i czasowniki ruchu (\de{Ich bin gestanden}). Pełniejsze zestawienie: jednostka~35.

\section*{Kultura pracy}
\begin{center}\small
\begin{tabular}{@{}p{0.25\linewidth}p{0.67\linewidth}@{}}
\toprule
\textsf{\textbf{Obszar}} & \textsf{\textbf{Obserwacje (kulturowe, nie prawne)}}\\
\midrule
Punktualność & spotkania zaczynają się o czasie; spóźnienie wymaga uprzedzenia\\
Bezpośredniość & rzeczowa, ale uprzejma; unika się ostrej krytyki\\
Konsensus & decyzje po dyskusji; hierarchia mniej widoczna niż w Polsce\\
\de{du}/\de{Sie} & w IT często \de{du}; klient i urząd \de{Sie}\\
Prywatność & rozmowa o zarobkach i życiu osobistym jest powściągliwa\\
Spotkania towarzyskie & \de{Apéro}, \de{Feierabendbier}, lunch zespołu\\
\bottomrule
\end{tabular}
\end{center}
To uogólnienia; zespoły międzynarodowe w Zurychu potrafią być bardzo różne.

\section*{Szukanie pracy i umowa}
Źródła ofert: portale ogłoszeń, strony firm, LinkedIn, rekruterzy. Dokumenty: CV, list motywacyjny, świadectwa, referencje (\de{Dossier}). Rozmowa
bywa wieloetapowa: HR, zespół, zadanie techniczne, rozmowa z kierownikiem. Wynagrodzenie podaje się brutto rocznie w CHF; często \de{13.\ Monatslohn}
\verify{(zwyczajowe, nie ustawowe)}. Okres próbny i wypowiedzenie wynikają z umowy i z ustawy \verify{(art.\ 335b i 335c Obligationenrecht)}. Urlop minimum cztery tygodnie
rocznie \verify{(art.\ 329a OR)}. Szczegóły: jednostka~43.

\section*{Przeprowadzka i formalności}
Zameldowanie w gminie, ubezpieczenie zdrowotne (obowiązkowe), pozwolenie na pobyt, konto bankowe, telefon, podatki (czasem \de{Quellensteuer}), ubezpieczenia społeczne
(\de{AHV}, \de{Pensionskasse}). Wszystko z terminami zależnymi od kantonu \verify{(sprawdź na ch.ch i stronie gminy)}. Szczegóły: jednostka~35.

\section*{Zwroty na pierwszy miesiąc}
\begin{center}\small
\begin{tabular}{@{}p{0.43\linewidth}p{0.47\linewidth}@{}}
\toprule
\textsf{\textbf{Sytuacja}} & \textsf{\textbf{Zwrot}}\\
\midrule
Przedstawienie się & \de{Ich bin Piotr, ich arbeite im Backend-Team.}\\
Prośba o wolniej & \de{Könnten Sie bitte etwas langsamer sprechen?}\\
Prośba o standard & \de{Könnten Sie bitte Hochdeutsch sprechen?}\\
Potwierdzenie & \de{Ich habe das verstanden, ich fasse kurz zusammen \dots}\\
Brak zrozumienia & \de{Das habe ich nicht ganz verstanden.}\\
Podziękowanie & \de{Merci vielmal.} \de{Danke für Ihre Hilfe.}\\
\bottomrule
\end{tabular}
\end{center}
